\section{Empirical Characterization}
\label{sec:characterization}

The empirical characterization emphasizes benchmark structure rather than a large model bake-off. The release metadata establish the scale and heterogeneity of the benchmark, and the released candidate-row summaries add full label-distribution reporting without changing that focus.

Across 277,995,072 candidate rows, \texttt{y\_loss} has mean 7.7561, standard deviation 4.6542, and 5/25/50/75/95th percentiles of 2/4/7/11/16. The corresponding \texttt{y\_value} summary mirrors those values by sign, as expected from the benchmark definition. The resulting label surface is broad but still bounded by the finite horizon, which keeps both candidate-level regression errors and decision-level regret interpretable. Table~\ref{tab:decision-breakdown}'s exactly equal horizon counts (787,762 rows at each of horizons 4, 8, and 16) reflect the same underlying full-cache-miss decisions relabeled at each configured horizon, whereas its capacity counts correspond to genuinely different full-cache-miss events produced by replaying each trace at each cache size; pooled statistics in this paper therefore mix horizons and capacities and should be read as aggregate release summaries rather than per-configuration results.

The same result summaries also show that the simple linear-score selector often lands in the benchmark-optimal set even though its underlying value regression is weak. That combination is useful descriptive context for later benchmark users: the release is not trivial, but a lightweight heuristic already clears a reproducible lower bar. The evaluator's grouped decision count now matches the manifest-backed decision-view count exactly (Table~\ref{tab:label-selection-summary}); an earlier decision-key bug that double-counted a small number of decisions has been fixed and reconciled. Because the dataset is heavily tied at the decision level, a uniform-random selector's expected optimal-selection rate and regret are included alongside the linear-score selector's observed metrics so that headline numbers are not read in isolation from tie density.

\begin{table}[!h]
  \centering
  \footnotesize
  \caption{Decision-view-backed benchmark composition by family, capacity, horizon, and split.}
  \label{tab:decision-breakdown}
  \begin{tabularx}{\columnwidth}{@{}P{0.22\columnwidth}L@{}}
    \toprule
    dimension & counts \\
    \midrule
    trace\_family & wiki2018=598,560; alibaba-block=576,219;\newline
    metakv=454,053; twemcache=397,494;\newline
    metacdn=336,960 \\
    capacity & 32=613,314; 64=599,172;\newline
    128=582,678; 256=568,122 \\
    horizon & 4=787,762; 8=787,762; 16=787,762 \\
    split & test=242,418; train=1,746,882; val=373,986 \\
    \bottomrule
  \end{tabularx}
\end{table}


\begin{table}[t]
  \centering
  \footnotesize
  \caption{Decision-view-backed candidate-count statistics. These values summarize the number of candidates per decision in the preserved release.}
  \label{tab:candidate-count-stats}
  \begin{tabular}{@{}lrrrr@{}}
    \toprule
    group & min & med & mean & max \\
    \midrule
    overall & 32 & 64 & 117.63 & 256 \\
    alibaba-block & 32 & 64 & 118.35 & 256 \\
    metacdn & 32 & 64 & 117.12 & 256 \\
    metakv & 32 & 64 & 119.52 & 256 \\
    twemcache & 32 & 64 & 111.52 & 256 \\
    wiki2018 & 32 & 64 & 119.85 & 256 \\
    \bottomrule
  \end{tabular}
\end{table}


\begin{table}[t]
  \centering
  \footnotesize
  \caption{Candidate-label and best-candidate-selection summaries from validated result records.}
  \label{tab:label-selection-summary}
  \begin{tabularx}{\columnwidth}{@{}P{0.50\columnwidth}L@{}}
    \toprule
    statistic & value \\
    \midrule
    candidate rows summarized & 277,995,072 \\
    \texttt{y\_loss} mean / std & 7.7561 / 4.6542 \\
    \texttt{y\_loss} q05 / q25 / q50 / q75 / q95 & 2 / 4 / 7 / 11 / 16 \\
    \texttt{y\_value} mean / std & -7.7561 / 4.6542 \\
    \texttt{y\_value} q05 / q25 / q50 / q75 / q95 & -16 / -11 / -7 / -4 / -2 \\
    manifest-backed decision rows & 2,363,286 \\
    best-candidate evaluator decision count & 2,363,286 (exact match) \\
    \midrule
    \multicolumn{2}{@{}l}{\textit{all decisions (secondary, ecological view)}} \\
    best-candidate optimal-selection rate & 0.8753 \\
    best-candidate regret mean / median / max & 0.1251 / 0 / 3 \\
    random-selector expected optimal-selection rate & 0.9912 \\
    random-selector expected regret & 0.0088 \\
    \midrule
    \multicolumn{2}{@{}l}{\textit{discriminative decisions only (primary selector-quality view)}} \\
    discriminative decisions (opt.\ candidate count $<$ candidate count) & 764,282 (32.34\%) \\
    best-candidate optimal-selection rate & 0.6145 \\
    best-candidate mean regret & 0.3867 \\
    random-selector expected optimal-selection rate & 0.9729 \\
    random-selector expected regret & 0.0273 \\
    \bottomrule
  \end{tabularx}
\end{table}

